\section{Модуль I. Векторные пространства и евклидова геометрия}

\subsection{Базовый вопрос I.1: Дайте определение вещественного векторного пространства}
\textbf{Ответ:} Множество $V$ с операцией сложения $+ \colon V \times V \to V$ и умножения на скаляр $\cdot \colon
\mathbb{R} \times V \to V$, удовлетворяющее 8 аксиомам: 1) $u+v = v+u$; 2) $(u+v)+w = u+(v+w)$; 3) $\exists 0 \colon
v+0=v$; 4) $\exists (-v) \colon v+(-v)=0$; 5) $1 \cdot v = v$; 6) $\alpha(\beta v) = (\alpha\beta)v$; 7) $\alpha(u+v) =
\alpha u + \alpha v$; 8) $(\alpha+\beta)v = \alpha v + \beta v$.

\subsection{Базовый вопрос I.2: Что такое подпространство?}
\textbf{Ответ:} Непустое подмножество $U \subseteq V$, которое само является векторным пространством относительно тех же
операций.
Критерий: $0 \in U$ и $\forall u, v \in U, \alpha, \beta \in \mathbb{R} \implies \alpha u + \beta v \in U$.

\subsection{Базовый вопрос I.3: Что такое linear combination, span и generating set?}
\textbf{Ответ:} Линейная комбинация — выражение вида $\sum_{i=1}^k c_i v_i$.
Линейная оболочка $\spanop(S)$ — множество всех линейных комбинаций векторов из $S$.
Порождающая система — множество векторов $S$, для которого $\spanop(S) = V$.

\subsection{Базовый вопрос I.4: Определения linear independence, basis, dimension, coordinates}
\textbf{Ответ:}
\begin{itemize}
    \item \textbf{Линейная независимость:} $\sum c_i v_i = 0 \iff \text{все } c_i = 0$.
    \item \textbf{Базис:} линейно независимая порождающая система.
    \item \textbf{Размерность $\dim V$:} число векторов в базисе.
    \item \textbf{Координаты:} единственный набор скаляров $(x_1, \dots, x_n)$ такой, что $x = \sum x_i e_i$.
\end{itemize}

\subsection{Базовый вопрос I.5: Что такое inner product, norm, distance и orthogonality?}
\textbf{Ответ:} Скалярное произведение — симметричная, положительно определенная билинейная форма.
Норма — длина вектора $\|x\| = \sqrt{\langle x, x \rangle}$.
Расстояние — метрика $d(x, y) = \|x - y\|$.
Ортогональность — условие $\langle x, y \rangle = 0$.

\subsection{Базовый вопрос I.6: Что такое orthogonal projection?}
\textbf{Ответ:} Ортогональная проекция вектора $v$ на подпространство $U$ — это единственный вектор $p \in U$, такой что
вектор ошибки $v - p$ ортогонален подпространству $U$: $(v - p) \perp U$.

\subsection{Базовый вопрос I.7: Что такое hyperplane и normal vector?}
\textbf{Ответ:} Гиперплоскость в $\mathbb{R}^n$ — аффинное подпространство размерности $n - 1$, задаваемое уравнением
$\langle n, x \rangle = c$.
Вектор $n \ne 0$ называется вектором нормали (он перпендикулярен гиперплоскости).

\newpage

\section{Модуль II. Матрицы, линейные отображения, ранг и СЛУ}

\subsection{Базовый вопрос II.1: Что такое matrix multiplication?}
\textbf{Ответ:} Операция $C = AB$, где $c_{ij} = \sum_{k=1}^p a_{ik} b_{kj}$ (скалярное произведение $i$-й строки $A$ на
$j$-й столбец $B$).
Является координатной записью композиции операторов.

\subsection{Базовый вопрос II.2: Что такое linear transformation и как matrix его представляет?}
\textbf{Ответ:} Отображение $T \colon V \to W$, сохраняющее линейные комбинации: $T(\alpha u + \beta v) = \alpha T(u) +
\beta T(v)$.
Столбцы матрицы оператора представляют собой координаты образов базисных векторов.

\subsection{Базовый вопрос II.3: Определения inverse, determinant, rank и nullspace}
\textbf{Ответ:}
\begin{itemize}
    \item \textbf{Обратная матрица $A^{-1}$:} удовлетворяет $A A^{-1} = A^{-1} A = I$.
    \item \textbf{Определитель $\det A$:} знакопеременная полилинейная функция строк, коэффициент изменения
    ориентированного объёма.
    \item \textbf{Ранг $\rank A$:} число линейно независимых столбцов (размерность образа).
    \item \textbf{Ядро $\Null A$:} множество решений уравнения $Ax = 0$.
\end{itemize}

\subsection{Базовый вопрос II.4: Сформулируйте rank-nullity theorem}
\textbf{Ответ:} Для линейного отображения $T \colon V \to W$:
\[ \dim V = \rank(T) + \nullity(T). \]

\subsection{Базовый вопрос II.5: Как записывается и интерпретируется СЛУ?}
\textbf{Ответ:} Запись $Ax = b$.
Интерпретация строк (row picture): пересечение $m$ гиперплоскостей.
Интерпретация столбцов (column picture): представление вектора $b$ в виде линейной комбинации столбцов матрицы $A$.

\subsection{Базовый вопрос II.6: Что такое Gaussian elimination?}
\textbf{Ответ:} Алгоритм последовательного исключения неизвестных с помощью элементарных преобразований строк для
приведения матрицы к верхнетреугольному (ступенчатому) виду.

\subsection{Базовый вопрос II.7: Что такое LU decomposition?}
\textbf{Ответ:} Факторизация матрицы в виде $A = LU$, где $L$ — нижнетреугольная матрица с единицами на диагонали, а $U$
— верхнетреугольная матрица.

\newpage

\section{Модуль III. Собственные значения и матричные разложения}

\subsection{Базовый вопрос III.1: Дайте определение eigenvalue и eigenvector}
\textbf{Ответ:} Ненулевой вектор $v \ne 0$ и скаляр $\lambda$ такие, что $Av = \lambda v$.

\subsection{Базовый вопрос III.2: Как найти eigenvalues и eigenvectors?}
\textbf{Ответ:} 1) Найти корни характеристического уравнения $\det(A - \lambda I) = 0$.
2) Для каждого $\lambda$ найти ненулевые решения системы $(A - \lambda I)v = 0$.

\subsection{Базовый вопрос III.3: Что такое algebraic и geometric multiplicity?}
\textbf{Ответ:} Алгебраическая кратность — кратность корня в многочлене $\det(A - \lambda I) = 0$.
Геометрическая кратность — размерность ядра $\dim \ker(A - \lambda I)$.

\subsection{Базовый вопрос III.4: Что такое eigendecomposition и diagonalization?}
\textbf{Ответ:} Представление $A=P \Lambda P^{-1}$, где столбцы $P$ — собственные векторы, а диагональ $\Lambda$ —
собственные значения.
Возможно тогда и только тогда, когда существует базис из собственных векторов.

\subsection{Базовый вопрос III.5: Сформулируйте real spectral theorem}
\textbf{Ответ:} Всякая вещественная симметричная матрица $A = A^\top$ имеет вещественный спектр и ортогонально
диагонализируема: $A = Q \Lambda Q^\top$, где $Q$ — ортогональная матрица ($Q^\top Q = I$).

\subsection{Базовый вопрос III.6: Что такое orthogonal matrix?}
\textbf{Ответ:} Квадратная матрица $Q$, для которой $Q^\top Q = Q Q^\top = I$ ($Q^{-1} = Q^\top$).
Сохраняет длины векторов и углы ($\|Qx\| = \|x\|$).

\subsection{Базовый вопрос III.7: Что такое positive definite, QR и Cholesky?}
\textbf{Ответ:}
\begin{itemize}
    \item \textbf{Positive definite:} $x^\top A x > 0$ для всех $x \ne 0$ (все $\lambda_i > 0$).
    \item \textbf{QR-разложение:} $A = QR$ ($Q$ ортонормирована, $R$ верхнетреугольная).
    \item \textbf{Разложение Cholesky:} $A = L L^\top$ ($L$ нижнетреугольная с положительной диагональю).
\end{itemize}

\newpage

\section{Модуль IV. SVD и итоговые инструменты}

\subsection{Базовый вопрос IV.1: Что такое SVD?}
\textbf{Ответ:} Сингулярное разложение произвольной матрицы: $A = U \Sigma V^\top$, где $U$ и $V$ — ортогональные
матрицы, а $\Sigma$ содержит неотрицательные сингулярные числа на диагонали.

\subsection{Базовый вопрос IV.2: Как singular values связаны с eigenvalues?}
\textbf{Ответ:} Сингулярные числа матрицы $A$ равны квадратным корням из собственных значений матрицы $A^\top A$:
$\sigma_i(A) = \sqrt{\lambda_i(A^\top A)}$.
Для положительно полуопределенной матрицы сингулярные числа совпадают с собственными значениями.

\subsection{Базовый вопрос IV.3: Что такое reduced SVD?}
\textbf{Ответ:} Компактная форма $A = U_r \Sigma_r V_r^\top$, где отброшены нулевые сингулярные числа и соответствующие
столбцы матриц $U$ и $V$.
Размерность $\Sigma_r$ равна $r \times r$, где $r = \rank A$.

\subsection{Базовый вопрос IV.4: Каков geometric meaning of SVD?}
\textbf{Ответ:} Отображение единичной сферы в эллипсоид в три этапа: поворот в пространстве прообразов ($V^\top$),
масштабирование по осям на величины $\sigma_i$ ($\Sigma$), и поворот полученного эллипсоида в пространстве образов
($U$).

\subsection{Базовый вопрос IV.5: Что такое low-rank approximation и pseudoinverse?}
\textbf{Ответ:}
\begin{itemize}
    \item \textbf{Low-rank approximation:} усеченная сумма $A_k = \sum_{i=1}^k \sigma_i u_i v_i^\top$ ($k < r$),
    наилучшее приближение матрицы $A$ рангом $k$.
    \item \textbf{Pseudoinverse (Мура-Пенроуза):} матрица $A^+ = V \Sigma^+ U^\top$, обеспечивающая обобщенное
    решение задачи наименьших квадратов $x = A^+ b$ с минимальной евклидовой нормой.
\end{itemize}